\documentclass[../main.tex]{subfiles}

\graphicspath{{\subfix{../images/}}}

\begin{document}


\clearpage
\thispagestyle{empty}
\vspace*{7cm}
\section{Operational amplifiers}
\vspace*{1.5cm}

\subsection{Introductory analysis}

\subsubsection{Ideal operational amplifiers}


The simplest approach to studying circuits based on operational amplifiers involves using a \textit{"black-box"} model (which does not specify the internal structure of the device in question). 

The most commonly used representation for the operational amplifier is that of a \textit{"triangle"} which is equipped with two input terminals, two power supply terminals (often omitted in circuits), and one output terminal; the input terminals, marked with the symbols "$+$" and "$-$" (also referred to as \textit{"non-inverting input"} and \textit{"inverting input"} respectively), are the inputs for the signals that the operational amplifier is intended to amplify; the power supply terminals, as the name suggests, are intended to bias and supply power to the circuit contained within the operational amplifier, in order to ensure its proper operation. 

The equations governing the operation of an \textit{"ideal"} operational amplifier are:

\begin{equation}
\begin{cases}
	i_+ = i_- \rightarrow 0 \\
	v_d = v_+ - v_{-} \rightarrow 0
\end{cases}
\end{equation}


These equations are fundamental to the study of a general circuit containing one (or more) operational amplifiers. 

Since the op-amp has a gain tending to infinity, it can be inferred that, in order to have a finite output—that is, so that the result of the product of the differential voltage $v_d$ (voltage between the "$+$" and "$-$" terminals) and the op-amp’s differential gain Ad is finite $v_d$ must tend to 0. 

Consequently, in the ideal operational amplifier, the voltage drop across the terminals is nearly zero and the input current is nearly zero, regardless of the differential resistance present between the input terminals. 

To simplify the calculations, one can assume that the operational amplifier terminals present a differential resistance $r_d \rightarrow \infty$ to the input currents. 
In summary, the fundamental characteristics of the ideal operational amplifier are:

\begin{enumerate}
	\item[\textbf{-}] Differential gain approaching infinity \\
	\item[\textbf{-}] Differential input resistance approaching infinity (a convenient but not necessary assumption)
	\item[\textbf{-}]Output resistance approaching 0
	\item[\textbf{-}]Differential input voltage approaching 0
	\item[\textbf{-}]Input currents approaching 0
\end{enumerate}

Let's try applying what we've just learned to a practical example (\autoref{ex1}). We consider the following circuit:

\begin{figure}[h]
\begin{center}
	
\begin{tikzpicture}[transform shape]
	\node[ground] at (5, 6.98){};
	\node[ground] at (7, 4.375){};
	\draw (5, 6.98) to[american resistor, l={$R_1$}] (7, 6.98);
	\draw (7, 8.5) to[american resistor, l={$R_2$}] (10, 8.5);
	\draw (7, 4.375) to[european voltage source, v=$v_{in}$] (7, 6);
	\node[op amp] at (8.815, 6.49){};
	\draw (7, 8.5) -| (7, 6.98);
	\draw (7, 6) -- (7.625, 6);
	\draw (7.625, 6.98) |- (7, 6.98);
	\draw (10, 8.5) -| (10, 6.5);
	\draw (10.005, 6.49) -- (11.005, 6.49);
	\node[circ] at (7, 6.98){};
	\node[circ] at (10.005, 6.49){};
	\node[ocirc](N1) at (11.005, 6.49){} node[anchor=south] at (N1.north){$v_{out}$};
\end{tikzpicture}

\label{ex1}
\caption{Non-inverting amplifier.}	
\end{center}	
\end{figure}

This circuit, as we will see shortly, is a non-inverting amplifier, that is, it amplifies a signal without inverting its phase (or shifting it by 180°). 
As an amplifier, it will have a certain gain, defined as the ratio of the output voltage, $v_{out}$, to the input voltage, $v_{in}$. 
It can be easily seen, taking into account the operating equations of the
device, that:

\begin{equation}
	v_{out} \cdot \frac{R_1}{R_2 + R_1} = v_-
\end{equation}

Since we can say that:

\begin{equation}
	\frac{v_{out}}{v_{in}} = \frac{R_1 + R_2}{R_1} = \left( 1 + \frac{R_2}{R_1} \right)
\end{equation}

At this point, we would like to draw some conclusions regarding the practical example just presented:

\begin{enumerate}[label=\textbf{-}]
	\item In this first part of the discussion, the operational amplifier will be used in a feedback configuration. As with any type of system that incorporates feedback, negative feedback will result in the effects already familiar from introductory electronics courses: changes in input or output impedance, an increase in bandwidth, and more.
	\item When the feedback is connected to the "$-$" terminal of the operational amplifier, it is \textit{"negative"} since the signal always opposes the input, reducing it. Feedback on the non-inverting terminal will be positive.
	\item In automatic control theory, feedback systems are often modeled with an amplification block $A$ and a feedback block $\beta$ (\autoref{ex2}). In the case of operational amplifiers, it is often easy to distinguish the A block from the $\beta$ block; the $\beta$ block is the circuit (passive network, in this case) capable of \textit{"feeding back"} a portion of the output signal to the input. Since, with this topology, the signal \textit{"fed back to the inverting terminal"} is equal to:
	
	\begin{equation}
		v_{out} \cdot \frac{R_1}{R_1 + R_2} = v_- = v_+
	\end{equation}
	
	we can also say that:
	
	\begin{equation}
		\beta = \frac{R_1}{R_1 + R_2}
	\end{equation}
	
	
	\begin{figure}[h]
	\begin{center}
	
	\begin{tikzpicture}[transform shape]
	\node[shape=rectangle, draw, line width=1pt, minimum width=0.965cm, minimum height=0.965cm] at (5.5, 9.5){};
	\node[shape=rectangle, minimum width=0.965cm, minimum height=0.965cm] at (5.5, 9.5){} node[anchor=center, align=center, text width=0.577cm, inner sep=6pt] at (5.5, 9.5){$A$};
	\node[shape=rectangle, draw, line width=1pt, minimum width=0.965cm, minimum height=0.965cm] at (5.5, 8){} node[anchor=center, align=center, text width=0.577cm, inner sep=6pt] at (5.5, 8){$\beta$};
	\node[shape=circle, draw, line width=1pt, minimum width=0.465cm] at (3.75, 9.5){};
	\draw[line width=0.8pt] (4, 9.5) -- (5, 9.5);
	\draw[line width=0.8pt] (3.75, 8) -- (3.75, 8) -- (5, 8);
	\draw[line width=0.8pt] (6, 9.5) -- (7, 9.5);
	\node[shape=rectangle, minimum width=0.465cm, minimum height=0.465cm] at (3.125, 9.875){} node[anchor=center, align=center, text width=0.077cm, inner sep=6pt] at (3.125, 9.875){$+$};
	\draw[line width=0.8pt, -latex] (2.625, 9.5) -- (3.5, 9.5);
	\draw[line width=0.8pt, -latex] (7, 9.5) -- (7.875, 9.5);
	\draw[line width=0.8pt] (7, 9.5) -- (7, 8);
	\draw[line width=0.8pt, -latex] (7, 8) -- (6.125, 8) -- (6, 8);
	\draw[line width=0.8pt, -latex] (3.75, 8) -- (3.75, 9.25);
	\node[shape=rectangle, minimum width=0.59cm, minimum height=0.527cm] at (4.062, 8.906){} node[anchor=center, align=center, text width=0.202cm, inner sep=6pt] at (4.062, 8.906){$-$};
	\draw (2.625, 7.5) to[open, v=$v_{in}$, voltage/bump b=0] (2.625, 9.5);
	\draw (7.375, 7.5) to[open, v=$v_{out}$, voltage/bump b=0] (7.375, 9.5);
	\end{tikzpicture}
	
	\label{ex2}	
	\caption{Example of feedback system.}
	\end{center}	
	\end{figure}
	
	
\end{enumerate}

\subsubsection{Non-ideal operational amplifiers}

Our discussion involves several approximations: real operational amplifiers have characteristics that deviate from the ideal model.

We have seen that a non-inverting amplifier can be constructed simply by selecting the resistors in the feedback loop to achieve a certain gain. But is the gain really the only thing that matters? To put the question another way: does using 1 $\Omega$ and 9 $\Omega$ resistors produce the same result as using 1 M$\Omega$ and 9 M$\Omega$ resistors, or 1 m$\Omega$ and 9 m$\Omega$ resistors? The answer is obviously no: real operational amplifiers exhibit non-ideal effects such that they are affected by the order of magnitude of the resistors used.  As can be seen from a study of the operational amplifier at the transistor level, we will see why it is not possible to use just any resistor. Essentially, the non-idealities are:

\begin{enumerate}[label=\textbf{-}]
	\item The gain $A_d$ is not infinite
	\item Non-infinite differential resistance rd and non-zero output resistance;
	\item Non-zero input currents;
	\item Non-zero differential voltage;
	\item Non-infinite input voltage and output voltage/current dynamics.
\end{enumerate}


Let us proceed step by step, presenting models that are progressively more refined than the ideal one; note that the approach used here does not explain the non-idealities, but rather studies the circuit’s behavior in the presence of a non-ideality. 
The justification for the presence of non-idealities, starting with an examination of the internal workings of the operational amplifier, will be provided later in this discussion.

\subsubsection*{Differential gain}

We present an initial refinement of our model: of the non-idealities listed above, let us consider the fact that $A_d < \infty$. The fact that $A_d$ is not infinite implies that, in order to have a non-zero output, a $v_d \neq 0$ is required. The new model of the device, therefore, will be that shown in figure \ref{ex3}.


\begin{figure}[h]
\begin{center}

\begin{tikzpicture}[scale = 0.75, transform shape]
	\node[ground] at (1, 11){};
	\node[ground] at (3, 7){};
	\node[ground] at (8, 8){};
	\draw (1, 11) to[american resistor, l={$R_1$}] (3, 11);
	\draw (7.125, 15) to[american resistor, l={$R_2$}] (9.125, 15);
	\draw (3, 7) to[european voltage source, v=$v_{in}$] (3, 9);
	\draw (3, 11) -- (5, 11);
	\draw (3, 9) -- (5, 9);
	\node[ocirc] at (5, 11){};
	\node[ocirc] at (5, 9){};
	\draw (5, 11) to[open, v^>=$v_d$, voltage/distance from node=0.8, voltage/bump b=0.2] (5, 9);
	\draw[line width=0.8pt, dash pattern={on 3.2pt off 3.2pt}] (3.75, 14) -| (3.75, 6);
	\draw[line width=0.8pt, dash pattern={on 3.2pt off 3.2pt}] (3.75, 6) -- (12, 10) -- (3.75, 14);
	\draw (8, 10) -- (13, 10);
	\draw (8, 8) to[european controlled voltage source, v^>=$A_dv_d$, voltage/distance from node=1, voltage=straight] (8, 10);
	\draw (3, 11) -| (3, 15) -- (7.125, 15);
	\draw (9.125, 15) -| (12.5, 10);
	\node[circ] at (12.5, 10){};
	\node[circ] at (3, 11){};
	\node[ocirc](N1) at (13, 10){} node[anchor=north] at (N1.text){$v_{out}$};
	\node[shape=rectangle, minimum width=0.59cm, minimum height=0.59cm] at (4.063, 9.313){} node[anchor=center, align=center, text width=0.202cm, inner sep=6pt] at (4.063, 9.313){$+$};
	\node[shape=rectangle, minimum width=0.59cm, minimum height=0.59cm] at (4.125, 11.313){} node[anchor=center, align=center, text width=0.202cm, inner sep=6pt] at (4.125, 11.313){$-$};
\end{tikzpicture}

\caption{First circuit model of the non-ideal operational amplifier: $A_d < \infty$.}
\label{ex3}	
\end{center}	
\end{figure}

We will then have that:

\begin{equation}
	v_- = v_{in} - v_d = v_{out} \cdot \beta
\end{equation}

However, one could also say that:

\begin{equation}
	v_d = \frac{v_{out}}{A_d}
\end{equation}

From here we get that:

\begin{equation}
	v_i - \frac{v_{out}}{A_d} = v_{out} \cdot \beta \rightarrow v_{out} \left( \beta + \frac{1}{A_d} \right)
\end{equation}

Therefore:

\begin{equation}
	v_{out} \cdot \frac{\beta A_d + 1}{A_d} = v_{in} \rightarrow \frac{v_{out}}{v_{in}} = \frac{A_d}{1 + \beta A_d} = \frac{1}{\beta} \cdot \frac{\beta A_d}{1 + \beta A_d} = \frac{1}{\beta} \cdot \frac{T}{1 + T}
\end{equation}


In feedback circuit theory, $T \triangleq \beta A_d $ is the \textit{"loop gain"}. Note, from this model, that in practical cases the deviation of the circuit’s behavior from the ideal case is very small: to achieve a 50\% deviation from the ideal case, the loop gain, $T$ would need to be 1. In reality, even the worst operational amplifiers might have a differential gain, $A_d$ of 10,000; by imposing a very high gain value on the circuit, $\beta$ could be around 1/1000. Few real circuits require a single amplification stage to have such a high gain, because this would cause problems with the circuit’s frequency response. In any case:

\begin{equation}
	T \simeq \frac{10000}{1000} = 10
\end{equation}

Even under these extremely harsh conditions, there is still a 10\% deviation between the gain of the actual circuit and that of the ideal circuit, which is generally acceptable.

\subsubsection*{Input impedance}

We can complicate our discussion by introducing other non-idealities: the impedances of operational amplifiers. Let’s consider a differential resistance $r_d$ that is not infinite (we will not consider the output resistance for now, so the voltage is still taken from an ideal voltage source). We want to calculate this, and to do so, we introduce a known test voltage source, $V_x$, at the input. In order to determine the input impedance, we calculate the current flowing out of $V_x$;

\begin{equation}
	V_x = I_x \cdot r_d + R_1 \cdot (I_{out} + I_x)
\end{equation}

We obtain that:

\begin{equation}
	v_d = r_d \cdot I_x \hspace{1cm} v_{out} = A_dv_d = A_dr_dI_x
\end{equation}

Furthermore:

\begin{equation}
	v_{out} = R_2 I_{out} + R_1(I_{out} + I_x) \rightarrow A_dr_dI_x = R_2I_{out} + R_1(I_{out} + I_x)
\end{equation}

By factoring out the term $I_{out}$:

\begin{equation}
	I_{out}(R_1 + R_2) = A_dr_dI_x - R_1Ix \rightarrow I_{out} = \frac{A_dr_dI_x - R_1I_x}{R_1 + R_2}	
\end{equation}

Substituting this into the expression for $V_x$, we can determine:

\begin{equation}
	V_x = I_xr_d + R_1I_x + \left( \frac{A_dr_dI_x - R_1I_x}{R_1 + R_2} \right) R_1
\end{equation}

From which:

%% TYPO TO FIX I GUESS!

\begin{equation}
	V_x = I_xr_d + \frac{R_1}{R_1 + R_2} A_dr_dI_x + \frac{R_1R_2 + R_1^2 - R_2^2}{R_1 + R_2} I_x
\end{equation}

Recalling that $\beta = \frac{R_1}{R_1 + R_2}$ we obtain that:

\begin{equation}
	z_{in} = \frac{V_x}{I_x} = r_d (1+ \beta A_d) + \beta R_2
\end{equation}

The second term can often be considered negligible compared to the first
(it contributes to increasing the impedance, so neglecting it yields a
\textit{"worst-case"} scenario); what is interesting is that even this model, which is significantly more refined than the ideal one, does not entail any particular changes to the circuit’s behavior: the series-comparison feedback significantly increases the circuit’s input impedance, making the assumption of an ideal operational amplifier acceptable once again in many of our calculations.\\

We end this second case by giving the previous seen model with the addition of the input resistance $r_d$ and by showing the currents $I_{out}$ and $I_x$ (\autoref{ex4}).

\begin{figure}[h]
\begin{center}

\begin{tikzpicture}[scale = 0.75, transform shape]
	\node[ground] at (1, 11){};
	\node[ground] at (3, 7){};
	\node[ground] at (8, 8){};
	\draw (1, 11) to[american resistor, l={$R_1$}] (3, 11);
	\draw (7.125, 15) to[american resistor, l={$R_2$}, i^<=$I_{out}$, current/distance=1] (10.5, 15);
	\draw (3, 7) to[european voltage source, v=$V_{x}$, i^>=$I_x$, current/distance=0.6] (3, 9);
	\draw (3, 11) -- (5, 11);
	\draw (3, 9) -- (5, 9);
	\draw[line width=0.8pt, dash pattern={on 3.2pt off 3.2pt}] (3.75, 14) -| (3.75, 6);
	\draw[line width=0.8pt, dash pattern={on 3.2pt off 3.2pt}] (3.75, 6) -- (12, 10) -- (3.75, 14);
	\draw (8, 10) -- (13, 10);
	\draw (3, 11) -| (3, 15) -- (7.125, 15);
	\draw (10.5, 15) -| (12.5, 10);
	\node[circ] at (12.5, 10){};
	\node[circ] at (3, 11){};
	\node[ocirc](N1) at (13, 10){} node[anchor=north] at (N1.text){$v_{out}$};
	\draw (8, 10) to[european controlled voltage source, v_<=$A_dv_d$] (8, 8);
	\draw (5, 11) to[american resistor, l_={$r_d$}, label distance=0.14cm, v^=$v_d$, voltage/bump b=0.5, voltage/shift=1.6] (5, 9);
	\node[shape=rectangle, minimum width=0.715cm, minimum height=0.715cm] at (4.125, 9.375){} node[anchor=center, align=center, text width=0.327cm, inner sep=6pt] at (4.125, 9.375){$+$};
	\node[shape=rectangle, minimum width=0.715cm, minimum height=0.715cm] at (4.125, 11.375){} node[anchor=center, align=center, text width=0.327cm, inner sep=6pt] at (4.125, 11.375){$-$};
\end{tikzpicture}

\label{ex4}
\caption{Second circuit diagram of the operational amplifier:
input resistor.}
\end{center}
\end{figure}

\subsubsection*{Exit impedance}

In order to refine the model already presented, we must consider the possible effects of the output impedance. Let us therefore consider the operational amplifier model shown in \autoref{ex5}.

To calculate the output impedance, we connect a test voltage generator to it—the usual $V_x$—and thus consider all other independent generators in the circuit to be turned off (the driven ones, of course, are not!). The current $I_x$ will consist of two components: one entering the branch of the driven generator $A_dv_d$ and one entering the branch of $R_2$; we can simplify the analysis by noting that $I_2 \ll I_1$:  since $r_o$ is a much smaller resistance than $R_1$, $R_2$, and even their parallel combination, we could say that $I_x  \simeq I_1$. Neglecting $I_2$ results in a higher impedance value than the actual one, which is acceptable when we want to verify that the output impedance is reasonably low. 

We therefore have that:

\begin{equation}
	I_x \simeq I_1 = \frac{V_x - A_dv_d}{r_o}
\end{equation}

However, we also know that $v_d$ can be expressed as:

\begin{equation}
	v_d = - \beta V_x = -\frac{R_1}{R_1 + R_2} V_x
\end{equation}

We can then say that:


\begin{equation}
		I_x \simeq \frac{V_x + A_d \beta V_x}{r_o}
\end{equation}


From here we get:

\begin{equation}
	\frac{I_x}{V_x} \simeq \frac{1+ \beta A_d}{r_o} \hspace{2cm} \beta A_d = T
\end{equation}

Therefore:

\begin{equation}
	Z_o = \frac{V_x}{I_x} \simeq \frac{r_o}{1 + T}
\end{equation}



Suppose we have a resistance of 100 $\Omega$, which is higher than what is found in most real-world amplifiers; if the loop gain were around 1000, we would reduce the resistance by three orders of magnitude, bringing it down to 100 m$\Omega$! 
We can therefore say that this circuit (non-inverting amplifier) is a good voltage amplifier: high input impedance and low output impedance.

\begin{figure}[h]
\begin{center}

\begin{tikzpicture}[scale = 0.75, transform shape]
	\node[ground] at (1, 11){};
	\node[ground] at (3, 7){};
	\node[ground] at (8, 8){};
	\draw (1, 11) to[american resistor, l={$R_1$}] (3, 11);
	\draw (7.125, 15) to[american resistor, l={$R_2$}, i^<=$I_{2}$, current/distance=1] (10.5, 15);
	\draw (3, 11) -- (5, 11);
	\draw (3, 9) -- (5, 9);
	\draw[line width=0.8pt, dash pattern={on 3.2pt off 3.2pt}] (3.75, 14) -| (3.75, 6);
	\draw[line width=0.8pt, dash pattern={on 3.2pt off 3.2pt}] (3.75, 6) -- (12, 10) -- (3.75, 14);
	\draw (3, 11) -| (3, 15) -- (7.125, 15);
	\draw (10.5, 15) -| (12.5, 10);
	\node[circ] at (12.5, 10){};
	\node[circ] at (3, 11){};
	\draw (8, 10) to[european controlled voltage source, v_<=$A_dv_d$] (8, 8);
	\draw (5, 11) to[american resistor, l_={$r_d$}, label distance=0.14cm, v^=$v_d$, voltage/bump b=0.5, voltage/shift=1.6] (5, 9);
	\node[shape=rectangle, minimum width=0.715cm, minimum height=0.715cm] at (4.125, 9.375){} node[anchor=center, align=center, text width=0.327cm, inner sep=6pt] at (4.125, 9.375){$+$};
	\node[shape=rectangle, minimum width=0.715cm, minimum height=0.715cm] at (4.125, 11.375){} node[anchor=center, align=center, text width=0.327cm, inner sep=6pt] at (4.125, 11.375){$-$};
	\draw (3, 7) -| (3, 9);
	\draw (8, 10) to[american resistor, l={$r_o$}, i^<=$I_1$, current/distance=1] (10, 10);
	\draw (10, 10) -- (13, 10);
	\draw (13, 10) to[european voltage source, v<=$V_x$, i=$I_x$] (13, 8);
	\node[ground] at (13, 8){};
\end{tikzpicture}

\caption{Operational amplifier model with non-zero output impedance.}
\label{ex5}	
\end{center}	
\end{figure}


\subsubsection{Voltage follower}



\subsubsection{Transresistance amplifier}
\subsubsection{Inverting amplifier}
\subsubsection{Current amplifier}	
\subsubsection{Transconductance amplifier}

%\subsection{Current Mirrors}
%\subsubsection{BJT Current Mirror}
%\subsection{MOSFET Current Mirror}
%\subsection{BJT Differential Stage}
%
%\subsection{Differential stage}
%\subsubsection{Differential mode and common-mode}
%\subsubsection{BJT differential stage}
%\subsubsection{MOSFET differential stage}
%
%\subsection{Simplified schematic of an operational amplifier}
%
%\subsection{Bipolar transistor power stages}
%\subsubsection{Bipolar power transistors}
%\subsubsection{Power classes}
%\subsubsection{Class A}
%\subsubsection{Class B}
%\subsubsection{Class AB}
%\subsubsection{Design Example}
%
%\subsection{Real operational amplifier models}
%\subsubsection{Voltage and current offset}
%\subsubsection{Common-mode input dynamics}
%\subsubsection{Output dynamics}
%\subsubsection{Input impedances}
%\subsubsection{Differential gain}
%\subsubsection{CMOS operational amplifier}
%
%\subsection{Amplifier sizing}
%\subsubsection{Specifications}
%\subsubsection{Project}
%
%\subsection{Frequency response}
%\subsubsection{Transfer function}
%\subsubsection{Loop gain}
%\subsubsection{Bandwidth-gain product}
%\subsubsection{Slew rate}









\end{document}